\subsection{Sequence-to-Sequence Learning}

After introducing the main categories of sequential problems, we now focus on a particularly important family: \definition{Sequence-to-Sequence Learning} (often abbreviated as \textbf{Seq2Seq}). In this setting, the goal is to learn a mapping from an \textbf{input sequence} to an \textbf{output sequence}. Unlike standard classification, where the output is a single label, here the model must produce a structured sequence of predictions, one element at a time.

\highspace
This is a natural framework for many real-world applications where both the input and the output carry sequential structure. Some examples are:
\begin{itemize}
    \item \textbf{Image Captioning}: the input is a single image, and the output is a variable-length sentence describing it.
    \item \textbf{Sentiment Analysis}: the input is a variable-length sentence or tweet, and the output is a single label such as positive or negative.
    \item \textbf{Machine Translation}: the input is a sentence in one language, and the output is a sentence in another language.
\end{itemize}
Although these examples do not all have the same input/output shape, they share the same core principle: the \hl{model must learn to generate outputs \textbf{conditioned on the input}} and, when the \hl{output is sequential} (i.e., a sentence or a sequence of labels), it \hl{must do so \textbf{autoregressively}}, i.e. one token at a time (predicting the next token based on the previous ones).

\highspace
\textcolor{Green3}{\faIcon[regular]{lightbulb} \textbf{Main Idea.}} In sequence-to-sequence learning, we do not want to predict just one target value, but an entire target sequence:
\begin{equation*}
    y_1, y_2, \ldots, y_n
\end{equation*}
Given an input sequence:
\begin{equation*}
    x_1, x_2, \ldots, x_m
\end{equation*}
Where, in general, the two lengths \emph{can} be different:
\begin{equation*}
    m \neq n \quad \text{but not necessarily}
\end{equation*}
This is what makes Seq2Seq models more general than aligned many-to-many models, where the output is forced to follow the same temporal structure as the input.

\highspace
Although several types of sequential problems exist, \hl{sequence-to-sequence learning is studied in greater detail because it is the most general and expressive setting.} It captures the core difficulties of sequence modeling: handling variable-length structured inputs, generating variable-length outputs autoregressively, and learning a conditional distribution $P\left(y \mid x\right)$. Simpler sequential tasks, such as many-to-one classification or aligned sequence labeling, can often be interpreted as restricted special cases of this broader framework.